\subsection{PHÂN LOẠI VÀ PHƯƠNG PHÁP GIẢI TOÁN}
\begin{dang}{Tính trung vị của mẫu số liệu ghép nhóm}
\end{dang}
\begin{vd}%[1K3B9-2]
	Thời gian (phút) truy cập internet mỗi buổi tối của một số học sinh được cho trong bảng sau:
	\begin{center}
		\begin{tabular}{|c|c|c|c|c|c|c|}
			\hline
			Thời gian (phút)	& $\left[9{,}5;12{,}5 \right)$ & $\left[12{,}5;15{,}5 \right)$ & $\left[15{,}5;18{,}5 \right)$ & $\left[18{,}5;21{,}5 \right)$ & $\left[21{,}5;24{,}5 \right)$ \\
			\hline
			Số học sinh&$3$	& $12$ & $15$ &$24$ & $2$ \\
			\hline
		\end{tabular}	
	\end{center}
	Tính trung vị của mẫu số liệu ghép nhóm này.
	\loigiai{
		Cỡ mẫu là $n=3+12+15+24+2=56$.\\
		Gọi $x_1,\,\ldots,\,x_{56}$ là thời gian vào internet của $56$ học sinh và giả sử dãy này đã được sắp xếp theo thứ tự tăng dần. Khi đó, trung vị là $\dfrac{x_{28}+x_{29}}{2}$. Do $2$ giá trị $x_{28},\,x_{29}$ thuộc nhóm $\left[15{,}5;18{,}5 \right)$ nên nhóm này chứa trung vị. Do đó, $p=3$; $a_3=15{,}5$; $m_3=15$; $m_1+m_2=3+12=15$; $a_4-a_3=3$ và ta có $$M_e=15{,}5+\dfrac{\dfrac{56}{2}-15}{15}\cdot 3=18{,}1.$$
	}
\end{vd}\dongcham{5}

\begin{vd}%[1T5B2-1]
	Kết quả khảo sát cân nặng của $ 25 $ quả bơ ở một lô hàng cho trong bảng sau:
	\begin{center}
		\begin{tabular}{|c|c|c|c|c|c|}
			\hline 
			Cân nặng  (g)	& $ \left[150 ; 155 \right) $ & $ \left[ 155 ; 160\right) $ & $ \left[160 ; 165\right) $ & $ \left[ 165 ; 170\right) $ & $ \left[170 ; 175 \right) $ \\ 
			\hline 
			Số quả bơ	& $ 1 $ & $ 7 $ & $ 12 $ & $ 3 $ & $ 2 $ \\ 
			\hline 
		\end{tabular} 
	\end{center}
	Hãy tìm trung vị của mẫu số liệu ghép nhóm trên.
	\loigiai{
		Gọi $ x_1; x_2; \ldots ; x_{25} $ là cân nặng của $ 25$ quả bơ xếp theo thứ tự không giảm.\\
		Do $ x_1\in \left[150 ; 155 \right) $; $ x_2, \ldots, x_8 \in \left[ 155 ; 160\right) $; $ x_9, \ldots, x_{20} \in \left[ 160 ; 165\right) $ nên trung vị của mẫu số liệu $ x_1; x_2; \ldots; x_{25} $ là $ x_{13}\in \left[ 160 ; 165\right)$.\\
		Ta xác định được $ n=25 $, $p=3 $, $ a_3=160$, $ m_3=12$, $m_1+m_2=1+7=8$, $a_4-a_3=165-160=5$.\\
		Vậy trung vị của mẫu số liệu ghép nhóm là $$ M_e=160+\dfrac{\dfrac{25}{2}-8}{12}\cdot 5 =161{,}875.$$
	}
\end{vd}\dongcham{4}

\begin{vd}%[1D5H2-2]%[Dự án đề kiểm tra Toán 11 HKI NH24-25-Võ Thị Thùy Trang]%[Lê Quý Đôn - Tp HCM]
	Trong một hội thao, thời gian chạy $200$ m của một nhóm các vận động viên được ghi lại ở bảng sau
	\begin{center}
		\begin{tabular}{|c|c|c|c|c|c|}
			\hline
			Thời gian (giây) &$[21;21{,}5)$& $[21{,}5;22)$ &$[22;22{,}5)$&$[22{,}5;23)$ &$[23;23{,}5)$\\
			\hline
			Số vận động viên & $5$ & $12$ & $32$ & $45$ & $30$ \\
			\hline
		\end{tabular}
	\end{center}
	Dựa vào bảng số liệu trên, ban tổ chức muốn chọn ra khoảng $50\%$ số vận động viên chạy nhanh nhất để tiếp tục thi vòng $2$. Ban tổ chức nên chọn các vận động viên có thời gian chạy không quá bao nhiêu giây (kết quả làm tròn đến hàng phần chục)?
	\loigiai{Trung vị của bảng số liệu trên là $M_\text{e}=u_m+\dfrac{\dfrac{n}{2}-C}{n_m}\left(u_{m+1}-u_m\right)$.\\
		$M_\text{e}=22{,}5+\dfrac{\dfrac{124}{2}-49}{45}\left(23-22{,}5\right)=\dfrac{1019}{45}\approx 22{,}64$ (giây).}
\end{vd}
\dongcham{4}

\begin{dang}{Tìm tứ phân vị của mẫu số liệu ghép nhóm}
\end{dang}

\begin{vd}%[1K3B9-3]
	Thời gian (phút) truy cập internet mỗi buổi tối của một số học sinh được cho trong bảng sau:
	\begin{center}
		\begin{tabular}{|c|c|c|c|c|c|c|}
			\hline
			Thời gian (phút)	& $\left[9{,}5;12{,}5 \right)$ & $\left[12{,}5;15{,}5 \right)$ & $\left[15{,}5;18{,}5 \right)$ & $\left[18{,}5;21{,}5 \right)$ & $\left[21{,}5;24{,}5 \right)$ \\
			\hline
			Số học sinh&$3$	& $12$ & $15$ &$24$ & $2$ \\
			\hline
		\end{tabular}	
	\end{center}
	Tìm tứ phân vị thứ nhất $Q_1$ và tứ phân vị thứ ba $Q_3$ của mẫu số liệu ghép nhóm.
	\loigiai{
		Cỡ mẫu là $n=3+12+15+24+2=56$.\\
		Tứ phân vị thứ nhất $Q_1$ là $\dfrac{x_{14}+x_{15}}{2}$. Do $2$ giá trị $x_{14},\,x_{15}$ thuộc nhóm $\left[12{,}5;15{,}5 \right)$ nên nhóm này chứa $Q_1$. Do đó, $p=2$; $a_2=12{,}5$; $m_2=12$; $m_1=3$; $a_3-a_2=3$ và ta có $$Q_1=12{,}5+\dfrac{\dfrac{56}{4}-3}{12}\cdot 3=15{,}25.$$
		Với tứ phân vị thứ ba $Q_3$ là $\dfrac{x_{42}+x_{43}}{2}$. Do $2$ giá trị $x_{42},\,x_{43}$ thuộc nhóm $\left[18{,}5;21{,}5 \right)$ nên nhóm này chứa $Q_3$. Do đó, $p=4$; $a_4=18{,}5$; $m_4=24$; $m_1+m_2+m_3=3+12+15=30$; $a_5-a_4=3$ và ta có $$Q_3=18{,}5+\dfrac{\dfrac{3\cdot 56}{4}-30}{24}\cdot 3=20.$$
	}
\end{vd}\dongcham{6}
\begin{vd}%[1C5B1-4]
	Bảng bên cho biết tần số ghép nhóm số liệu thống kê cân nặng của $40$ học sinh lớp $11A$ trong một trường trung học phổ thông (đơn vị: kilôgam). Xác định tứ phân vị của mẫu số liệu ghép nhóm.
		\begin{center}
			\begin{tabular}{|c|c|c|c|c|c|c|}
					\hline 
					Nhóm &
					$\left[30;40\right)$ &
					$\left[40;50\right)$ &
					$\left[50;60\right)$ &
					$\left[60;70\right)$ &
					$\left[70;80\right)$ &
					$\left[80;90\right)$ \\
					\hline 
					Tần số &
					$2$ &
					$10$ &
					$16$ &
					$8$ &
					$2$ &
					$2$ \\
					\hline
			\end{tabular}
		\end{center}
	\loigiai{
		Cỡ mẫu là $n=40$.
		\begin{itemize}
			\item 
			Tứ phân vị thứ nhất $Q_1$ là $\dfrac{x_{10}+x_{11}}{2}$. Do $2$ giá trị $x_{10},\,x_{11}$ thuộc nhóm $\left[40;50 \right)$ nên nhóm này chứa $Q_1$. Do đó, $p=2$; $a_2=40$; $m_2=10$; $m_1=2$; $a_3-a_2=10$ và ta có $$Q_1=40+\dfrac{\dfrac{40}{4}-2}{10}\cdot 10=48~(\mathrm{(kg)}).$$
			%%
			\item Với tứ phân vị thứ hai $Q_2$ là $\dfrac{x_{20}+x_{21}}{2}$. Do $2$ giá trị $x_{20},\,x_{21}$ thuộc nhóm $\left[50; 60\right)$ nên nhóm này chứa $Q_2$. Do đó, $p=3$; $a_3=50$; $m_3=16$; $m_1+m_2=12$; $a_4-a_3=60-50=10$ và ta có $$Q_2=50+\dfrac{\dfrac{ 40}{2}-12}{16}\cdot 10=55~(\mathrm{(kg)}).$$	
			%%
			\item
			Với tứ phân vị thứ ba $Q_3$ là $\dfrac{x_{30}+x_{31}}{2}$. Do $2$ giá trị $x_{30},\,x_{31}$ thuộc nhóm $\left[60; 70\right)$ nên nhóm này chứa $Q_3$. Do đó, $p=4$; $a_4=60$; $m_4=8$; $m_1+m_2+m_3=28$; $a_5-a_4=70-60=10$ và ta có $$Q_3=60+\dfrac{\dfrac{3\cdot 40}{4}-28}{8}\cdot 10=62{,}5~(\mathrm{(kg)}).$$	
		\end{itemize}
		Vậy tứ phân vị của mẫu số liệu trên là $48$, $55$ và $62{,}5$.
	}
\end{vd}\dongcham{6}
\begin{vd}%[1T5K2-2]
	Thời gian luyện tập trong một ngày (tính theo giờ) của một số vận động viên được ghi lại ở bảng sau:
	\begin{center}
		\begin{tabular}{|c|c|c|c|c|c|}
			\hline 
			Thời gian luyện tập (giờ)	& $ \left[ 0 ; 2\right) $ & $ \left[ 2 ; 4\right) $ & $ \left[ 4 ; 6\right) $ & $ \left[ 6 ; 8\right) $ & $ \left[8 ; 10 \right) $ \\ 
			\hline 
			Số vận động viên	& $ 3 $ & $ 8 $ & $ 12 $ & $ 12 $ & $ 4 $ \\ 
			\hline 
		\end{tabular} 
	\end{center}
\begin{tasks}
	\task Hãy xác định các tứ phân vị của mẫu số liệu đã cho.
	\task Huấn luyện viên muốn xác định nhóm gồm $ 25\% $ các vận động viên có số giờ luyện tập cao nhất. Hỏi huấn luyện viên nên chọn các vận động viên có thời gian luyện tập từ bao nhiêu giờ trở lên vào nhóm này?
\end{tasks}
	\loigiai{
		\begin{enumerate}[a)]
			\item Số vận động viên được khảo sát là $ n=3+8+12+12+4=39$.\\
			Gọi $ x_1 $; $ x_2 $; \ldots ;$ x_{39} $ là thời gian luyện tập của $ 39 $ vận động viên được xếp theo thứ tự không giảm.\\
			Ta có $ x_1, x_2, x_3 \in \left[ 0 ; 2\right) $; $ x_4, \ldots, x_{11}\in \left[2 ; 4 \right) $; $ x_{12}, \ldots, x_{23}\in \left[ 4 ; 6\right) $; $ x_{24}, \ldots, x_{35}\in \left[6;8\right) $; $ x_{36},\ldots , x_{39}\in \left[ 8 ; 10\right) $. Do đó đối với dãy số liệu $ x_1 $; $ x_2 $; \ldots; $ x_{39} $ thì
			\begin{itemize}
				\item Tứ phân vị thứ nhất là $ x_{10} $ thuộc nhóm $\left[ 2 ; 4\right)$;
				\item Tứ phân vị thứ hai là $ x_{20} $ thuộc nhóm $ \left[ 4 ; 6\right)$;
				\item Tứ phân vị thứ ba là $ x_{30} $ thuộc nhóm $ \left[ 6 ; 8\right) $.
			\end{itemize}
			Ta nói nhóm $ \left[2 ; 4 \right)$ là nhóm chứa tứ phân vị thứ nhất; nhóm $ \left[ 4 ; 6\right) $ là nhóm tứ phân vị thứ hai; nhóm $ \left[ 6 ; 8\right) $ là nhóm chứa tứ phân vị thứ ba.\\
			Tứ phân vị thứ hai của dãy số liệu $ x_1; x_2; \ldots ; x_{39}$ là $ x_{20}\in \left[ 4 ; 6\right) $. Do đó tứ phân vị thứ hai của mẫu số liệu ghép nhóm là $$Q_2=4+\dfrac{\dfrac{2\cdot 39}{4}-(3+8)}{12}\cdot(6-4)=\dfrac{65}{12}\approx 5{,}417.$$ 
			Tứ phân vị thứ nhất của dãy số liệu $ x_1; x_2; \ldots ; x_{39}$ là $ x_{10}\in \left[ 2 ; 4\right) $. Do đó tứ phân vị thứ nhất của mẫu số liệu ghép nhóm là $$Q_1=2+\dfrac{\dfrac{1\cdot 39}{4}-3}{8}\cdot(4-2)=\dfrac{59}{16}\approx 3{,}6875.$$ 
			Tứ phân vị thứ ba của dãy số liệu $ x_1; x_2; \ldots ; x_{39}$ là $ x_{30}\in \left[ 6 ; 8\right) $. Do đó tứ phân vị thứ ba của mẫu số liệu ghép nhóm là $$Q_3=6+\dfrac{\dfrac{3\cdot 39}{4}-(3+8+12)}{12}\cdot(8-6)=\dfrac{169}{24}\approx 7{,}042.$$ 
			\item Huấn luyện viên muốn xác định nhóm gồm $ 25\% $ các vận động viên có số giờ luyện tập cao nhất thì huấn luyện viên nên chọn các vận động viên có thời gian tập luyện từ $ Q_3 \approx 7{,}042$ (giờ).
		\end{enumerate}
		
		
	}
\end{vd}\dongcham{7}

\begin{vd}%[1T5K2-2]
	Một phòng khám thống kê số bệnh nhân đến khám bệnh mỗi ngày trong tháng $ 4 $ năm $ 2022 $ ở bảng sau:
	\begin{center}
		\begin{tabular}{|c|c|c|c|c|c|}
			\hline 
			\textbf{Số bệnh nhân}	& $ \left[1 ; 10 \right] $ & $ \left[ 11 ; 20\right] $ & $ \left[ 21 ; 30\right] $ & $ \left[ 31 ; 40\right] $ & $ \left[41 ; 50 \right] $ \\ 
			\hline 
			\textbf{Số ngày}	& $ 7 $ & $ 8 $ & $ 7 $ & $ 6 $ & $ 2 $ \\ 
			\hline 
		\end{tabular} 
	\end{center}
	\begin{enumerate}[a)]
		\item Hãy ước lượng các tứ phân vị của mẫu số liệu ghép nhóm trên.
		\item Quản lý phòng khám cho rằng có khoảng $ 25\% $ số ngày khám có nhiều hơn $ 35 $ bệnh nhân đến khám. Nhận định trên có hợp lý không?
	\end{enumerate}
	\loigiai{
		\begin{enumerate}[a)]
			\item Do số bệnh nhân đến khám là số nguyên nên ta hiệu chỉnh lại như sau:
			\begin{center}
				\begin{tabular}{|c|c|c|c|c|c|}
					\hline 
					\textbf{Số bệnh nhân}	& $ \left[0{,}5 ; 10{,}5 \right) $ & $ \left[ 10{,}5 ; 20{,}5\right) $ & $ \left[ 20{,}5 ; 30{,}5\right) $ & $ \left[ 30{,}5 ; 40{,}5\right) $ & $ \left[40{,}5 ; 50{,}5 \right) $ \\ 
					\hline 
					\textbf{Số ngày}	& $ 7 $ & $ 8 $ & $ 7 $ & $ 6 $ & $ 2 $ \\ 
					\hline 
				\end{tabular} 
			\end{center}
			Gọi $ x_1; x_2; \ldots; x_{30} $ là mẫu số liệu được xếp theo thứ tự không giảm.\\
			Ta có $ x_1, \ldots, x_{7}\in \left[ 0{,}5 ; 10{,}5\right) $; $ x_8, \ldots, x_{15} \in \left[ 10{,}5 ; 20{,}5\right) $, $ x_{16}, \ldots, x_{22} \in \left[ 20{,}5 ; 30{,}5\right) $; $ x_{23}, \ldots, x_{28} \in \left[ 30{,}5 ; 40{,}5\right) $; $ x_{29}, x_{30} \in \left[ 40{,}5 ; 50{,}5\right) $.\\
			Tứ phân vị thứ nhất của dãy số liệu $ x_1; x_2; \ldots; x_{30} $ là $ x_8 \in \left[10{,}5 ; 20{,}5 \right) $ nên tứ phân vị thứ nhất của mẫu số liệu ghép nhóm là $$ Q_1=10{,}5+\dfrac{\dfrac{30}{4}-7}{8}\cdot (20{,}5-10{,}5)=\dfrac{89}{8}\approx 11{,}125.$$
			Tứ phân vị thứ hai của dãy số liệu $ x_1; x_2; \ldots; x_{30} $ là $ \dfrac{1}{2}\left( x_{15}+x_{16}\right) $. Do $ x_{15} \in \left[ 10{,}5 ; 20{,}5 \right) $ và $ x_{16}\in \left[ 20{,}5 ; 30{,}5\right) $ nên tứ phân vị thứ hai của mẫu số liệu ghép nhóm là $ Q_2=20{,}5 $.\\
			Tứ phân vị thứ hai của dãy số liệu $ x_1; x_2; \ldots; x_{30} $ là $ x_{23}\in \left[ 30{,}5 ; 40{,}5\right)$ nên tứ phân vị thứ ba của mẫu số liệu ghép nhóm là $$Q_3= 30{,}5+\dfrac{\dfrac{3\cdot 30}{4}-(7+8+7)}{6}\cdot(40{,}5-30{,}5)=\dfrac{94}{3}\approx 31{,}333.$$
			\item Do $ Q_1\approx 11{,}125 $ nên nhận định $ 25\% $ số ngày khám nhiều hơn $ 35 $ bệnh nhân là không hợp lý.
		\end{enumerate}
	}
\end{vd}\dongcham{7}

\subsection{BÀI TẬP TỰ LUYỆN}


\begin{bt}%[1C5B1-5]
	\immini
	{
		Bảng bên cho ta bảng tần số ghép nhóm số liệu thống kê chiều cao của $40$ mẫu cây ở một vườn thực vật (đơn vị: centimét). Xác định trung vị, tứ phân vị của mẫu số liệu ghép nhóm trên.
	}
	{
		\begin{tabular}{|c|c|c|}
			\hline
			\textbf{Nhóm} & \textbf{Tần số}\\ 
			\hline
			$\left[30;40\right)$ & $4$ \\
			$\left[40;50\right)$ & $10$ \\
			$\left[50;60\right)$ & $14$\\
			$\left[60;70\right)$ & $6$\\
			$\left[70;80\right)$ & $4$ \\
			$\left[80;90\right)$ & $2$\\
			\hline
			& $n = 40$ \\
			\hline
		\end{tabular}
	}
	\loigiai{
			\begin{itemize}
				\item Tứ phân vị thứ 2 thuộc nhóm $3$ là nhóm $\left[50;60\right)$. Do đó, tứ phân vị thứ hai (cũng là trung vị) là
				\[
				Q_2 = M_e = 50 + \dfrac{20 - 14}{14} \cdot 10 = 54{,}3\text{ (cm)}.
				\]
				\item Tứ phân vị thứ nhất thuộc nhóm $2$ là nhóm $\left[40;50\right)$. Do đó, tứ phân vị thứ nhất là
				\[
				Q_1 = 40 + \dfrac{10 - 4}{10} \cdot 10 = 46\text{ (cm)}.
				\]
				\item Tứ phân vị thứ 3 thuộc nhóm $4$ là nhóm $\left[60;70\right)$. Do đó, tứ phân vị thứ ba là
				\[
				Q_3 = 60 + \dfrac{30 - 28}{6} \cdot 10 \approx 63{,}3\text{ (cm)}.
				\]
			\end{itemize}
	}
\end{bt}\cham{15}



\begin{bt}%[1C5B1-5]
	Mẫu số liệu dưới đây ghi lại tốc độ của $40$ ô tô khi đi qua một trạm đo tốc độ (đơn vị: km/h):
	\[
	\begin{array}{cccccccccc}
		48{,}5 & 43 & 50 & 55 & 45 & 60 & 53 & 55,5 & 44 & 65 \\
		51 & 62,5 & 41 & 44,5 & 57 & 57 & 68 & 49 & 46{,}5 & 53{,}5 \\
		61 & 49{,}5 & 54 & 62 & 59 & 56 & 47 & 50 & 60 & 61 \\
		49{,}5 & 52{,}5 & 57 & 47 & 60 & 55 & 45 & 47,5 & 48 & 61{,}5
	\end{array}
	\]
	\begin{enumerate}
		\item Lập bảng tần số ghép nhóm cho mẫu số liệu trên có sáu nhóm ứng với sáu nửa khoảng:
		\[
		[40 ; 45),[45 ; 50),[50 ; 55),[55 ; 60),[60 ; 65),[65 ; 70).
		\]
		\item Xác định số trung vị, tứ phân vị của mẫu số liệu ghép nhóm trên.
	\end{enumerate}
	\loigiai{
		\begin{enumerate}
			\item Ta có bảng tần số ghép nhóm của mẫu số liệu trên như sau:
			\begin{center}
				\begin{tabular}{|c|c|c|c|}
					\hline
					\textbf{Nhóm} & \textbf{Giá trị đại diện} & \textbf{Tần số} & \textbf{Tần số tích luỹ}\\ 
					\hline
					$\left[40;45\right)$ & $42{,}5$ & $4$ & $4$\\
					$\left[45;50\right)$ & $47{,}5$ & $11$ & $15$\\
					$\left[50;55\right)$ & $52{,}5$ & $7$ & $22$\\
					$\left[55;60\right)$ & $57{,}5$ & $8$ & $30$\\
					$\left[60;65\right)$ & $62{,}5$ & $8$ & $38$\\
					$\left[65;70\right)$ & $67{,}5$ & $2$ & $40$\\
					\hline
					& & $n = 40$ &\\
					\hline
				\end{tabular}
			\end{center}
			\item 
			\begin{itemize}
				\item Ta có trung vị (tứ phân vị thứ 2) thuộc nhóm nhóm $3$ là nhóm $\left[50;55\right)$. Do đó, tứ phân vị thứ hai (cũng là trung vị) là
				\[
				Q_2 = M_e = 50 + \dfrac{20-15}{7} \cdot 5 = 53{,}6\text{ (km/h)}.
				\]
				\item Ta có tứ phân vị thứ nhất thuộc nhóm $2$ là nhóm $\left[45;50\right)$. Do đó, tứ phân vị thứ nhất là
				\[
				Q_1 = 45 + \dfrac{10-4}{11} \cdot 5 \approx 47{,}7\text{ (km/h)}.
				\]
				\item Ta có tứ phân vị thứ 3 thuộc nhóm $4$ là nhóm $\left[55;60\right)$. Do đó, tứ phân vị thứ ba là
				\[
				Q_3 = 55 + \dfrac{30-22}{8} \cdot 5 = 60\text{ (km/h)}.
				\]
			\end{itemize}
		\end{enumerate}
	}
\end{bt}\cham{20}


\begin{bt}%[1T5K2-1]
	Trong tuần lễ bảo vệ môi trường, các học sinh khối $ 11 $ tiến hành thu nhặt vỏ chai nhựa để tái chế. Nhà trường thống kê kết quả thu nhặt vỏ chai của học sinh khối $ 11 $ ở bảng sau:
	\begin{center}
		\begin{tabular}{|c|c|c|c|c|c|}
			\hline 
			\textbf{Số vỏ chai nhựa}	& $ \left[ 11 ; 15\right] $ & $ \left[ 16 ; 29\right] $ & $ \left[21 ; 25 \right] $ & $ \left[ 26 ; 30\right] $ & $ \left[31 ; 35 \right] $ \\ 
			\hline 
			\textbf{Số học sinh}	& $ 53 $ & $ 82 $ & $ 48 $ & $ 39 $ & $ 18 $ \\ 
			\hline 
		\end{tabular} 
	\end{center}
	Hãy tìm trung vị của mẫu số liệu ghép nhóm trên.
	\loigiai{
		Do số vỏ chai là số nguyên nên ta hiệu chỉnh lại như sau:
		\begin{center}
			\begin{tabular}{|c|c|c|c|c|c|}
				\hline 
				\textbf{Số vỏ chai nhựa}	& $ \left[ 10{,}5 ; 15{,}5\right) $ & $ \left[ 15{,}5 ; 20{,}5\right) $ & $ \left[ 20{,}5 ; 25{,}5\right) $ & $ \left[ 25{,}5 ; 30{,}5\right) $ & $ \left[30{,}5 ; 35{,}5 \right) $ \\ 
				\hline 
				\textbf{Số học sinh}& $ 53 $ & $ 82 $ & $ 48 $ & $ 39 $ & $ 18 $ \\ 
				\hline 
			\end{tabular} 
		\end{center}
		Số học sinh tham gia thu nhặt vỏ chai nhựa là $$ n=53+82+48+39+18=240.$$
		Gọi $ x_1; x_2; \ldots ; x_{240} $ lần lượt là số vỏ chai $ 240 $ học sinh khối $ 11 $ thu nhặt được xếp theo thứ tự không giảm.\\
		Do $ x_1, \ldots, x_{53}\in \left[10{,}5 ; 15{,}5 \right) $; $ x_{54}, \ldots, x_{135}\in \left[ 15{,}5 ; 20{,}5\right)$ nên trung vị của mẫu số liệu $ x_1; x_2; \ldots;x_{240} $ là $$ \dfrac{1}{2}\left( x_{120}+x_{121}\right)\in \left[ 15{,}5 ; 20{,}5\right).$$
		Ta xác định được $ n=240$; $p=2$; $a_2=15{,}5$; $m_2=82$; $m_1=53$; $a_3-a_2=20{,}5-15{,}5=5$.\\
		Trung vị của mẫu số liệu ghép nhóm là $$ M_e=15{,}5+\dfrac{\dfrac{240}{2}-53}{82}\cdot 5=\dfrac{803}{41}\approx 19{,}59. $$
	}
\end{bt}\cham{20}
\begin{bt}%[1T5K2-1]
	Trong một hội thao, thời gian chạy $ 200 $ m của một nhóm các vận động viên được ghi lại ở bảng sau:
	\begin{center}
		\begin{tabular}{|c|c|c|c|c|c|}
			\hline 
			\textbf{Thời gian} \textbf{(giây)}& $ \left[21 ; 21{,}5 \right) $ & $ \left[ 21{,}5 ; 22\right) $ & $ \left[ 22 ; 22{,}5\right) $ & $ \left[ 22{,}5 ; 23\right) $ & $ \left[ 23 ; 23{,}5\right) $ \\ 
			\hline 
			\textbf{Số vận động viên} & $ 5 $ & $ 12 $ & $ 32 $ & $ 45 $ & $ 30 $ \\ 
			\hline 
		\end{tabular} 
	\end{center}
	Dựa vào bảng số liệu trên, ban tổ chức muốn chọn ra khoảng $ 50 \% $ số vận động viên chạy nhanh nhất để tiếp tục thi vòng $ 2 $. Ban tổ chức nên chọn các vận động viên có thời gian chạy không quá bao nhiêu giây?
	\loigiai{
		Số vận động viên tham gia là $$n=5+12+32+45+30=124.$$
		Gọi $ x_1; x_2; \ldots ; x_{124} $ lần lượt là thời gian chạy $ 200 $ m của $ 124 $ vận động viên được xếp theo thứ tự không giảm.\\
		Do $ x_1, \ldots, x_5 \in \left[ 21 ; 21{,}5 \right)$, $ x_6, \ldots, x_{17} \in \left[ 21{,}5 ; 22\right) $, $ x_{18}, \ldots, x_{49} \in \left[22 ; 22{,}5\right) $, $ x_{50},\ldots, x_{94} \in \left[ 22{,}5 ; 23\right) $ nên trung vị của mẫu số liệu $ x_1; x_2; \ldots ;x_{124} $ là
		$$\dfrac{1}{2}\cdot \left( x_{62}+x_{63}\right) \in \left[ 22{,}5 ; 23\right).$$
		Ta xác định được $ n=124 $; $p=4$; $a_4=22{,}5$; $m_4=45$; $m_1+m_2+m_3=5+12+32=49$; $a_5-a_4=23-22{,}5=0{,}5$.\\
		Trung vị của mẫu số liệu ghép nhóm là $$M_e=22{,}5 +\dfrac{\dfrac{124}{2}-49}{45}\cdot 0{,}5= \dfrac{1019}{45}\approx 22{,}64.$$
		Vậy ban tổ chức nên chọn các vận động viên có thời gian chạy không quá $ 22{,}64$ (giây) để tiếp tục thi vòng hai.
	}
\end{bt}\cham{16}

\begin{bt}%[1T5K2-2]
	Một hãng xe ô tô thống kê lại số lần gặp sự cố về động cơ của $ 100 $ chiếc xe cùng loại sau $ 2 $ năm sử dụng đầu tiên ở bảng sau:
	\begin{center}
		\begin{tabular}{|c|c|c|c|c|c|}
			\hline 
			\textbf{Số lần gặp sự cố}	& $ \left[ 1 ; 2\right] $ & $ \left[3 ; 4\right] $ & $ \left[ 5 ; 6\right] $ & $ \left[ 7 ; 8\right] $ & $ \left[ 9 ; 10\right] $ \\ 
			\hline 
			\textbf{Số xe}	& $ 17 $ & $ 33 $ & $ 25 $ & $ 20 $ & $ 5 $ \\ 
			\hline 
		\end{tabular} 
	\end{center}
	\begin{enumerate}[a)]
		\item Hãy ước lượng các tứ phân vị của mẫu số liệu ghép nhóm trên.
		\item Một người cho rằng có trên $ 25\% $ xe của hãng gặp không ít hơn $ 4 $ sự cố về đồng cơ trong $ 2 $ năm sử dụng đầu tiên. Nhận định trên có hợp lí không?
	\end{enumerate}
	\loigiai{
		\begin{enumerate}[a)]
			\item Do số lần gặp sự cố là số nguyên nên ta hiệu chỉnh lại như sau:
			\begin{center}
				\begin{tabular}{|c|c|c|c|c|c|}
					\hline 
					\textbf{Số lần gặp sự cố}	& $ \left[ 0{,}5 ; 2{,}5\right) $ & $ \left[ 2{,}5 ; 4{,}5\right) $ & $ \left[4{,}5 ; 6{,}5 \right) $ & $ \left[ 6{,}5 ; 8{,}5\right) $ &$ \left[8{,}5 ; 10{,}5 \right) $ \\ 
					\hline 
					\textbf{Số xe}	& $ 17 $ & $ 33 $ &$ 25 $ & $ 20 $ & $ 5 $ \\ 
					\hline 
				\end{tabular} 
			\end{center}
			Gọi $ x_1; x_2; \ldots; x_{100} $ là mẫu số liệu được xếp theo thứ tự không giảm.\\
			Ta có $ x_1, \ldots, x_{17}\in \left[ 0{,}5; 2{,}5\right)$; $ x_{18}, \ldots, x_{50}\in \left[ 2{,}5 ; 4{,}5\right)$; $ x_{51}, \ldots, x_{75}\in \left[ 4{,}5 ; 6{,}5\right)$; \\
			$ x_{76}, \ldots, x_{95}\in \left[ 6{,}5 ; 8{,}5\right)$; $ x_{96}, \ldots, x_{100}\in \left[ 8{,}5 ; 10{,}5\right) $.\\
			Tứ phân vị thứ hai của dãy số liệu $ x_1; x_2; \ldots; x_{100} $ là $ \dfrac{1}{2}\left( x_{50}+x_{51}\right) $. Do $ x_{50}\in \left[ 2{,}5 ; 4{,}5\right) $ và $ x_{51}\in \left[4{,}5 ; 6{,}5 \right) $ nên tứ phân vị thứ hai của mẫu số liệu ghép nhóm là $ Q_2=4{,}5$.\\
			Tứ phân vị thứ nhất của dãy số liệu $ x_1; x_2; \ldots; x_{100} $ là $ \dfrac{1}{2}\left( x_{25}+x_{26}\right) $. Do $ x_{25} $ và $ x_{26} $ thuộc nhóm $ \left[2{,}5 ; 4{,}5 \right)$ nên tứ phân vị thứ nhất của mẫu số liệu ghép nhóm là $$ Q_1=2{,}5+\dfrac{\dfrac{1\cdot100}{4}-17}{33}\cdot (4{,}5-2{,}5)=\dfrac{197}{66}\approx 2{,}98.$$
			Tứ phân vị thứ ba của dãy số liệu $ x_1; x_2; \ldots; x_{100} $ là $ \dfrac{1}{2}\left(x_{75}+x_{76} \right)$. Do $ x_{75}\in \left[4{,}5 ; 6{,}5 \right) $ và $ x_{76}\in \left[ 6{,}5 ; 8{,}5\right)$ nên tứ phân vị thứ ba của mẫu số liệu ghép nhóm là $ Q_3=6{,}5$.
			\item Do tứ phân vị thứ nhất $ Q_1\approx 2{,}98 $ nên nhận định trên là hợp lý.
		\end{enumerate}
	}
\end{bt}\cham{22}




